\documentclass[10pt]{article}
\usepackage[margin=20mm]{geometry}
\usepackage{amsmath,amssymb,mathtools,fontspec,unicode-math,xurl,hyperref}
\IfFontExistsTF{DejaVu Serif}{\setmainfont{DejaVu Serif}}{\setmainfont{Latin Modern Roman}}
\setmathfont{latinmodern-math.otf}
\setlength{\parindent}{0pt}\setlength{\parskip}{6pt}\sloppy
\title{Technical Verification of
Maximal Parahoric Arithmetic Transfers,
Resolutions and Modularity}
\author{AI-assisted technical review}
\date{}
\begin{document}
\maketitle
Abstract: The acquired 106-page author version is inspected through local transfer, modified intersections, maximal-order reduction, global modularity and double induction. Three failures are established in central intermediate claims: the unrestricted Cayley intersection identity, the maximal-order length formula and the exact Y-series constant-term identification. Explicit examples give \(1\ne0\), \(1\ne2\), and a nonzero constant discrepancy. These findings do not constitute counterexamples to the main arithmetic transfer conjecture: the independently evaluated geometric and analytic sides in the maximal-order example agree. Four additional defects have narrower scope. Two dependency questions remain Needs Evidence. The overall assessment is Major; scientific adjudication is pending.\par
\section*{Identity, version and scope}
The paper is Zhiyu Zhang, Maximal parahoric arithmetic transfers, resolutions and modularity, associated with Duke Mathematical Journal 174 (2025), 1–129, DOI 10.1215/00127094-2024-0023 (https://doi.org/10.1215/00127094-2024-0023). The exact reviewed file is arXiv:2112.11994v3 (https://arxiv.org/abs/2112.11994v3), 9 May 2024, with the internal date 10 May 2024 [paper]. It is the accepted author manuscript; equivalence to the publisher’s version has not been established. All page locations below are physical pages of this 106-page PDF, whose printed numbering agrees. Journal page numbers are not substituted.\par
The root covered pp. 1–13 and 90–106, with additional direct checks of the findings; the auxiliary assignments covered pp. 14–30, 31–60 and 61–89. The p. 61 end of Section 9 was also checked by the local-geometry role. Hash-bound image records distinguish actual inspections from rendering. The coverage ledger identifies external interfaces not independently reproved. The acquired author TeX was read as a source witness, not executed. No separate supplement or original computational repository was identified in bounded primary-source checks; absence of code is not a defect.\par
\section*{Assumptions and checks that survive}
For arithmetic transfer \(F/F_0\) is unramified quadratic and the residue characteristic is odd. A vertex lattice has \(L\subseteq L^\vee\subseteq\varpi^{-1}L\), rank \(n\), and type \(0\le t\le n\). The original hermitian space \(V\) and the special-homomorphism space \(\mathbb V\) have opposite determinant parity. For \(L=L^\flat\oplus O_Fe\), \(i=v((e,e))\) is zero or one and \(t(L^\flat)=t-i\). The intermediate-type main theorem additionally requires \(F_0/\mathbb Q_p\) unramified; the extreme types retain the broader odd-residue-characteristic scope. Part 1’s Jacquet–Rallis transfer includes \(p=2\), separately from the arithmetic assertion.\par
The global argument assumes \(F_0\ne\mathbb Q\), uses the specified CM data and regular projective integral model over \(O_E[\Delta^{-1}]\), and uses \(E=F\) for modification. Values lie in \[\mathbb R_\Delta
 =\mathbb R/\operatorname{span}_{\mathbb Q}
 \{\log\ell:\text{a place over }\ell\text{ is in }\Delta\}.\] The very-special-cycle restriction in Theorem 8.1 is retained for \(0<t<n\); general local modularity in Conjecture 8.2 is not claimed proved.\par
The product-node resolution has, after removal of Rees-chart torsion, charts \[x_1=x_0u,\quad y_0=uy_1,\quad x_0uy_1=\varpi,
 \qquad
 x_0=x_1v,\quad y_1=vy_0,\quad x_1vy_0=\varpi.\] These are regular strict semistable charts, with a projective-line exceptional fiber over the double linking locus. The graph lifts in Proposition 5.17 use a Cartier pullback of the center; they do not need the excessive generality criticized below. The unit-denominator Lemma 6.10 is preserved and used in Problem 1.\par
The Weil/Fourier-transfer interface in the paper’s reference [67], Theorem A.1, was checked in the 2021 typeset version [weil], physical pp. 107–109, printed pp. 969–971. Its local-field formulation imposes no \(p\ne2\) restriction and retains \(\eta(\det V)\). Its underlying local trace formula was not independently reproved. Finite-etale descent in Proposition 3.19 supplies the dyadic integral Hilbert 90 step. The separate \(p>2\) explanation does not justify a trace-zero integral basis at \(p=2\). Slips in \(\delta,\sqrt\delta\) and a sign in that alternative explanation are clarifications. The assigned checks also verify orthogonal-basis construction by residue trace, integral transitivity, volume-one stabilizer unfolding, and the two-root-group argument in the double uncertainty principle.\par
\section*{Central reduction and modularity findings}
Problem 1 (The given Cayley pair changes under the scalar twist).  Theorem 6.15(1), p. 45, asserts \(\widetilde{\operatorname{Int}}(G)=\widetilde{\operatorname{Int}}^{\mathcal Z}(c_U(G))\) for any relative regular semisimple \(G\), with the original untwisted pair. Its proof twists \(G\) to obtain a unit denominator, changing the right-hand side. Assessment: Major.\par
\subsection*{Every hypothesis and the explicit input}
Equation reference: eq:G, eq:cayley. Take \(F_0=\mathbb Q_3\), \(F=\mathbb Q_3(\iota)\), \(\iota^2=-1\). This extension is unramified. Hensel lifting gives \(r\in\mathbb Q_3\) with \(r^2=-2\), and \(a=r+\iota\) satisfies \(a\bar a=-1\). Take an original split rank-two space with self-dual \(L=L^\flat\oplus O_Fe\), \((e,e)=1\). Thus \(n=2,t=0,i=0\). In its nearby space use \(f,e_0\) with Gram matrix \(J=\operatorname{diag}(3,1)\). Put \[
 G=\begin{pmatrix}-2&a\\-3\bar a&-2\end{pmatrix}.\] Exact multiplication gives \(G^*JG=J\), \(\det G=1\). The vectors \(e_0,Ge_0\) are independent, with determinant \(-a\); \(T^2+4T+1\) has distinct roots. Relative regular semisimplicity holds. The reduced vector is nonzero, so its rank-one pair is also regular semisimple. Here \(d=-2,\ 1-d=3\). Formula (6.9), p. 43, gives \[
 c_U(G)=(-1,af/3),\qquad (af/3,af/3)=-1/3.\] By Proposition 5.7(3) this special divisor is empty, so the claimed right-hand side is zero.\par
\subsection*{An actual derived intersection of length one}
The moduli input exists: in the rank-one nearby space choose the special pair \(A=B=3^{-1}O_{\breve F}f\), with \(A^\vee=B^\vee=O_{\breve F}f\). It satisfies Proposition 7.6, and \(af\in B^\vee\) gives an integral special-fiber homomorphism by (7.14). Principal polarizations make its adjoint integral, with \(u^*u=uu^*=-3\). Thus [eq:G] (\#eq:G) is a genuine polarized automorphism of the framed product.\par
The central scalar \(-1\) acts trivially on Rapoport–Zink moduli: multiplication by \(-1\) on the object identifies the two framings. Hence \(\Gamma_G=\Gamma_{-G}\), whereas \[1-d(-G)=-1\in O_F^\times,\qquad c_U(-G)=(-1,af).\] Apply the correctly restricted Lemma 6.10 to \(-G\), without using the disputed theorem. For \(D=\delta(\mathcal N^\flat)=\mathcal Z(e_0)\), \[(D\times D)\cap\Gamma_G
 \simeq\mathcal Z(af)\subseteq\mathcal N^\flat=\mathcal N_1^{[0]}.\] Proposition 9.2, p. 56, gives \(\mathcal Z(af)=\operatorname{Spf}(O_{\breve F}/(3))\), since \(v(-3)=1\).\par
Equation reference: eq:onezero. At type zero the resolution is the identity. The ambient \(\mathcal N_2^{[0]}\) is a regular formal surface and \(D,GD\) are horizontal Cartier divisors. The first-factor graph pullback is Tor-independent. The remaining equation restricts to \(3\) times a unit on the DVR \(D\simeq\operatorname{Spf}O_{\breve F}\). It is a nonzerodivisor, so no higher Tor term cancels the ordinary length. Thus \[
 \widetilde{\operatorname{Int}}(G)=1,\qquad \widetilde{\operatorname{Int}}^{\mathcal Z}(c_U(G))=0.\] The rank-one canonical-lifting formula is the explicitly retained base-theory interface; a new proof of all canonical-lifting theory is not claimed.\par
\subsection*{Counterevidence, impact and correction}
The twist preserves the left side but changes the reduced norm from \(-1/3\) to \(-3\). The correct unit-denominator lemma is consistent with the example. One can impose \(1-d\in O_F^\times\), or explicitly choose the twist and redefine the pair afterward. This repairs the statement; every repaired general reduction has not been verified here. The identity is a central Section 6 bridge, cited in the boundary induction on p. 103. Theorem 6.15(3) may survive a correctly stated twisted reduction. No counterexample to that implication or the main arithmetic transfer identity is asserted.\par
Problem 2 (The maximal-order sum includes nonexistent fixed components).  Theorem 9.10, p. 59, sums over every algebra factor, omitting the signature/type admissibility in Theorem 9.6. This false formula is explicitly used to prove Theorem 9.14, p. 61. Assessment: Major.\par
Take the same \(F_0,F\), \(n=2,t=0\), and \[\mathbb V=\langle1,3\rangle,\qquad
 g=\operatorname{diag}(1,-1),\qquad u=(3,1).\] The first line is split and the second nonsplit. Since the eigenvalue difference \(2\) is a unit, \(O_F[g]=O_F\times O_F\) is an etale maximal order. Also \(\det[u,gu]=-6\ne0\). Both components are integral; no nonintegral inactive-component ambiguity is needed.\par
At total type zero every \(t_i\) is zero. Theorem 9.6 allows only \(i_0=2\): an active split line would require odd type. Thus \(\operatorname{Fix}(g)\) is one copy of \(\operatorname{Spf}O_{\breve F}\). The inactive \(u_1=3\) lifts everywhere between the canonical signature-\((0,1)\) factors. The active component has norm \(3\). Proposition 9.2 gives length one, and Corollary 9.7 identifies the modified derived fixed cycle with the ordinary smooth fixed scheme. Consequently \(\widetilde{\operatorname{Int}}^{\mathcal Z}(g,u)=1\). The printed formula instead gives \[\sum_{j=1}^{2}\max\{0,(v((u_j,u_j))+a_j)/2\}
 =\frac{2+0}{2}+\frac{1+1}{2}=2.\] The extra one comes from the nonexistent \(i_0=1\) component.\par
\subsection*{Independent analytic control}
Equation reference: eq:orb. A matching triple is \(\gamma=\operatorname{diag}(1,-1)\), \(b=(1,1)\), \(c=(9,3)\); its moments agree with those of \((g,u)\). The transfer factor is \(\eta(\det[b,\gamma b])=\eta(-2)=1\); the cyclic reference lattice is \(O_F^2\). Integral idempotents force stable lattices to be \(3^{s_1}O_F\oplus3^{s_2}O_F\). Support is exactly \(s_1\in\{-2,-1,0\},\ s_2\in\{-1,0\}\). With the volume-one stabilizer normalization, the six cosets give \[
 \operatorname{Orb}(z)=(3^{2z}-3^z+1)(1-3^z),\qquad
 \partial\operatorname{Orb}(0)=-\log3.\] This agrees with the actual geometric length one. The exact enumeration is separately recorded. Corollary 9.13 gives the same value: it retains the admissible active index, and the inactive rank-one orbital integral at zero is one.\par
\subsection*{Dependence and repair boundary}
Theorem 9.14 explicitly cites Theorem 9.10, Corollary 9.13 and the rank-one calculation. Those inputs do not agree as printed. Section 9 supplies a maximal-order input to global modification and induction. This establishes an invalid derivation tied to the core argument, although [eq:orb] (\#eq:orb) verifies the desired transfer identity in this example. A repaired sum must retain admissible components and inactive-factor integrality conditions. The complete general corrected formula is not established here. A single local field factor is not contradicted; global irreducibility alone does not force local irreducibility.\par
Problem 3 (The Y-series has the wrong prescribed constant).  Equation (13.5), p. 86, and Theorem 13.9, p. 89, identify the arithmetic series with a theta series for both test functions. For the Y-family the constant identification contradicts Theorem 8.7. Assessment: Major for the exact modularity claim.\par
\subsection*{The local discrepancy}
Equation reference: eq:ZY. Take \(n=2,t=1\) and a type-zero component \(C=\mathbb P^1_\Lambda\). It is very special by Definition 12.14. Equation (8.11), p. 54, gives, for \(u\ne0\), \[
 (\mathcal Z(u),C)=1_\Lambda(u),\qquad
 (\mathcal Y(u),C)=-q1_\Lambda(u).\] The unique locally constant extensions at zero have values \(1,-q\). The factor is deliberate and compatible with Theorem 8.1’s Fourier relation.\par
Equation reference: eq:constant. The global test functions have the same away-from-\(v_0\) part and local factors \(1_L,1_{L^\vee}\), both one at zero. Proposition 11.18, p. 73, identifies their nonzero coefficients with \(\mathcal Z,\mathcal Y\). Equation (13.5) prescribes the same constant \(-\phi_i(0)c_1(\widehat\omega)\) for both. Proposition 13.7 computes \(\deg(\omega|_C)=-1\). Retaining a common nonzero arithmetic degree factor \(\kappa\), the prescribed constants are both \(\kappa\); the theta constants are \(\kappa,-q\kappa\). For the arithmetic Fourier series \(F_Z,F_Y\), all nonconstant coefficients therefore cancel in \[
 F_Y+qF_Z=(q+1)\kappa\ne0.\] A nonzero constant is not a holomorphic modular form of positive parallel weight: a congruence element with nonzero lower-left entry has a nonconstant weight factor. Intersecting the two levels makes this the same modularity test. It rules out simultaneous exact modularity and the exact Y theta identity.\par
For a type-two component, (8.12) instead reads \[(\mathcal Z(u),C)=-q1_\Lambda(u),\qquad
 (\mathcal Y(u),C)=1_{\Lambda^\vee}(u).\] The zero limits are \(-q,1\). The two identities in Proposition 13.7 would force \(\deg\omega=q\) and \(q\deg\omega=1\), hence \(q^2=1\). Thus a uniform scalar change of the degree cannot repair that proposition on both kinds of component.\par
\subsection*{An admissible global instance}
The discrepancy can remain nonzero in \(\mathbb R_\Delta\). Take \[F_0=\mathbb Q(\sqrt5),\quad F=F_0(\iota),\quad
 V=\operatorname{diag}(1,11\sqrt5),\quad L=O_F^2.\] Choose the CM type sending \(\iota\) to \(+i\) at both real embeddings, with distinguished embedding \(\sqrt5<0\). Set \(\Delta=\{v:v\mid2\text{ or }v\mid5\}\), and choose sufficiently small level there. The signatures are \((1,1)\), \((2,0)\). The Galois CM field gives total reflex field \(E=F\). Ramification over \(2\) permits the paper’s nonempty CM datum. The prime \(11\) splits in \(F_0\), while \(-1\) is nonsquare modulo \(11\). At either place above \(11\), \(F_{0,v}=\mathbb Q_{11}\), the quadratic extension is unramified and \(L_v\) has type one. Propositions 12.15–12.19 provide basic projective-line components. Choose the common away-from-\(v_0\) Schwartz factor to have value one at zero. The resulting nonzero degree factor is a rational multiple of \(\log11\), outside the rational span of \(\log2,\log5\). Thus [eq:constant] (\#eq:constant) is nonzero in the required quotient.\par
\subsection*{Counterevidence and limited impact}
Theorem 13.5(1) asserts modularity only up to a constant; this example does not refute it. The vertical-annihilator version in Theorems 13.4 and 13.5(2) has a separate generic-fiber argument. Section 15 also uses nonzero coefficients of the genuinely defined theta series, and those coefficients remain correct. The bad constant alone does not disprove the final transfer identity. A repair must supply a suitable Y-family extension or vertical correction and reconcile both kinds of component. No complete global repair or author adoption is certified. The exact basic-locus theta/modularity theorem is a central claimed ingredient, so its established failure remains Major.\par
\section*{Additional defects with bounded local impact}
Problem 4 (A fixed Cayley parameter does not give every target norm).  Proposition 3.2 and Corollary 3.4, p. 18, and the reverse implication in Theorem 3.26, p. 25, omit an image constraint for fixed \(d\ne1\). Assessment: Minor.\par
For \(h=(e,e)\) and \(G=\left(\begin{smallmatrix}t&u\\w&d\end{smallmatrix}\right)\), unitarity requires \((u,u)=h(1-d\bar d)\). Thus arbitrary \(g,u,d\ne1\) cannot be reconstructed. For a relative-regular witness over the same unramified \(\mathbb Q_3\) extension, take \(V^\flat=\langle1\rangle,h=1,d=0\), and \((g,u_1)=(1,2)\). The target pair is regular semisimple, but its required last column \((2,0)\) has norm four, not one. Passing to orbits does not remove this invariant obstruction.\par
The corrected condition is, for \(\alpha=(u_1,u_1)/h,\ z=(1-d)^{-1}\), \[\operatorname{Tr}_{F/F_0}z=1+\alpha.\] If \(\alpha\) is integral, every finite-field trace fiber contains a nonzero element, including at \(q=2\). Surjectivity of the unramified integral trace lifts it to a unit \(z\) with the prescribed trace. This chooses \(d\) with unit denominator. Writing \(s=1-d\) and \(w_1(x)=(gx,u_1)/h\), the reconstruction is \[G=\begin{pmatrix}g-su_1w_1&su_1\\sw_1&d\end{pmatrix}.\] The trace condition makes its mixed pairings zero and its last-column norm \(h\), giving the correct inverse on the compatible locus. The supported lattice condition \(u_1\in\Lambda\cap h\Lambda^\vee\) makes \(\alpha\) integral. If \(\alpha\) is nonintegral, no conjugate can meet this lattice support, so the relevant matching integrals vanish. For integral \(\alpha\), choose the unit \(z\) as above. This variable-\(d\) and support split repairs the intended reverse transfer argument. A charitable reading of Corollary 3.4 as an isomorphism onto its image does not validate the explicit unrestricted inverse assertion. No central transfer conclusion is refuted by this local defect.\par
Problem 5 (Scalar normalization is stated for too many inputs and tests).  The paragraph after (3.3), p. 20, quantifies over every regular semisimple orbit and every Schwartz test function. Both require restrictions. Assessment: Minor.\par
For \(J_0=\operatorname{diag}(1,-1)\), let \[H=\begin{pmatrix}5/3&4/3\\4/3&5/3\end{pmatrix}.\] Then \(H^*J_0H=J_0,\det H=1\), and \(e_2,He_2\) are independent. Its \(d=5/3\) has valuation \(-1\). Every norm-one scalar \(\xi\) is a unit, so \(v(1-\xi d)=-1\). None makes the denominator a unit. Independently, scalar twisting does not preserve every Schwartz test: the invariant \(d\) separates the relative orbits of \(H,\xi H\) when \(d\ne\xi d\). A nonnegative compact locally constant bump near one orbit, supported in a small separated \(d\)-neighborhood, has a positive orbital integral there and zero on the other. In this example the conjugating \(U(1)\) is compact, so positivity and convergence present no difficulty.\par
The actual lattice-stabilizer characteristic functions are invariant under central norm-one units. A negative-valuation \(d\) is outside their support. One may first discard those zero contributions, then normalize the integral \(d\) cases and recompute the twisted Cayley pair. These restrictions restore the intended application; unrestricted Schwartz invariance is not a valid general theorem.\par
Problem 6 (Strict transform does not preserve every local dimension).  Proposition 5.14(2), p. 37, allows any closed formal subscheme \(Y\), without purity or a surviving-component condition. Assessment: Minor.\par
Let \(k\) be the residue field, \(X=\mathbb A_k^3,\ Z=V(z),\ Y=V(z)\cup V(x,y)\), with reduced ideal \((xz,yz)\). These are locally formally of finite type over \(O_{\breve F}\) through \(k\). Blowing up the ambient Cartier plane is the identity. The strict transform is the closure of \(Y\setminus Z\), the \(z\)-axis. It is nonempty, but at the origin its local dimension is one and \(\dim_0Y=2\). Regular/pure graph applications instead use Cartier pullback and Proposition 5.17; the checked Hecke applications use their smooth or etale projections. No actual construction is refuted. An appropriate surviving-pure-component condition is needed; a fully general replacement theorem is not certified.\par
Problem 7 (The single-incidence local model has excessive dimension).  The proof of Theorem 10.8, p. 66, retains only a line \(\ell\), a hyperplane \(H\), and \((\ell,H)=0\). This is not the claimed model. Assessment: Minor.\par
The text uses a vector space over \(W(k)[1/p]\) with test rings where \(p\) is nilpotent. Even correcting the scalar ring to an integral lattice leaves a substantive problem. For \(n=3,t=2\), reduction of \(\operatorname{diag}(p,p,1)\) has a two-dimensional radical \(R\). Every line in \(R\) and every hyperplane satisfy the sole condition, giving \[\mathbb P(R)\times(\mathbb P^2)^\vee
 \simeq\mathbb P^1\times\mathbb P^2\] of dimension three, versus claimed relative dimension two. At the stated linking point \(\ell=\langle e_1\rangle,H=R\) the same excess persists locally.\par
Proposition 5.9, explicitly cited as the proof’s model, keeps both polarized lattice maps \(\lambda,\lambda^\vee\), with composite \(\varpi\). Theorem 10.5 also supplies a separate semistability interface. These are counterevidence against treating the naive model as the actual integral model. Repair requires the correct linked-lattice deformation problem or justified flat closure and a local-ring calculation; changing only \(W(k)[1/p]\) to \(W(k)\) is insufficient. No counterexample to the actual semistability conclusion is established.\par
\section*{Double induction and unresolved interfaces}
\subsection*{Completing the relevant Lemma 15.1 argument}
Theorem 15.7 uses \(B=\{v_0\}\), \(k_1=-c,k_2=-c+1\), so \(k_1+k_2+2c=1\), and boundary valuations zero and minus one. The Fourier support argument makes the enlarged-unipotent difference \(f_i(hn(b))-f_i(h)\) constant. For this single-place application, Fourier uniqueness on the compact adelic unipotent quotient suffices; the printed absolute-convergence assertion for every continuous function is unnecessary. After all nonzero coefficients vanish, archimedean unipotent invariance and its conjugate by a rational lower unipotent in the level group give archimedean invariance. Strong approximation then gives constancy.\par
If the translation difference is \(c_b\), iteration gives \(mc_b\). Integral multiples of \(b\) lie in a compact \(p\)-adic set; continuity forces \(c_b=0\). Thus the original automorphic function, not merely its Whittaker coefficient, has enlarged right invariance. The dual Weyl relation gives the opposite root subgroup. The exponent inequality makes them generate the full local \(\mathrm{SL}_2\); an exact control is \[\begin{pmatrix}1&0\\1/3&1\end{pmatrix}
 \begin{pmatrix}1&-3\\0&1\end{pmatrix}
 \begin{pmatrix}1&0\\1/3&1\end{pmatrix}
 =\begin{pmatrix}0&-3\\1/3&0\end{pmatrix}.\] Strong approximation now applies legitimately to the original function. A positive-weight constant is zero. This completes the transition where the printed proof instead applies strong approximation directly to a Whittaker coefficient. No false Lemma 15.1 is inferred; the general multi-place external lemma was not fully reproved.\par
The cyclotomic example on p. 98 gives \(F_0=\mathbb Q\) for \(q=3,5\). A totally real extension split at \(p\) preserves the local field, e.g. \(\mathbb Q(\iota,\sqrt{13})/\mathbb Q(\sqrt{13})\) at \(3\), or \(\mathbb Q(\sqrt{-3},\sqrt{11})/\mathbb Q(\sqrt{11})\) at \(5\). This repairs the example-level \(F_0\ne\mathbb Q\) requirement and does not disprove globalization.\par
\subsection*{Needs Evidence: modification normalization}
Put \(A=\tau(Z^{\mathbb Q})[E:F]\). Equations (14.16), (15.5), (15.6), pp. 94,100, divide geometric and archimedean terms by \(A\), while (15.7) adds a bare theta sum. Theorem 15.6, p. 101, writes \(\Theta(\xi,\phi'_{C_\nu})=(\widehat Z^B(\xi,\phi_i),C_\nu)\) without that denominator. For \(T=\sum_\nu(\widehat Z^B(\xi,\phi_i),C_\nu)\), the auxiliary-cycle identity makes the displayed modified difference equal the original normalized difference plus \[(1-A^{-1})T.\] No explicit rescaling convention for these theta functions or cycles was found. The printed volume definition also suppresses the compact level present in (11.14). An unstated normalization may remove this residual; this is a precise question for reconciliation, not another established Major finding.\par
\subsection*{Needs Evidence: intermediate-type irreducibility}
Proposition 12.24, pp. 82–83, adapts external arguments for \(1<t<n-1\). The exact cited van Hoften version [vh], arXiv:2010.10496v1, was consulted at its setup and Theorem 3.7.1/Lemma 3.7.4. Besides properness and adjoint simplicity, that argument uses a quasi-split/tame setup and noncompact \(J_b\)-factors. Its auxiliary union of strata of dimension at most one is not literally the whole basic intersection named in the target proof. A connected subspace meeting every closure component can suffice, so a terminology difference alone is no counterexample. The target acknowledges quasi-split qualifications on p. 84.\par
The asserted adaptation, all group hypotheses and its two further external interfaces were not fully verified. This remains a bounded dependency limitation, not a claim that Proposition 12.24 is false. The \(n=2,t=1\) constant-term example uses the separate projective-line description, not this adaptation.\par
\section*{Overall impact and state}
Problems 1–3 establish central intermediate claims or derivations that are invalid in the reviewed version, independently of the unresolved normalization question. Problems 4–7 are local defects without demonstrated central damage. The explicit maximal-order example verifies the desired transfer identity using its actual geometric length. The Y-series finding leaves modularity up to a constant, the vertical-annihilator argument and nonzero theta coefficients as surviving statements. These distinctions do not certify fully repaired proofs of all main results.\par
Item | Assessment | Repair verification / remaining state\par
Problems 1–3 | Major | Explicit discrepancies checked; full general repairs unverified\par
Problems 4–7 | Minor | Local correction scope; actual main constructions not refuted\par
Two interfaces | Needs Evidence | Normalize modification; justify irreducibility adaptation\par
Scientific decision | Pending | Initial review, not independent acceptance\par
Author adoption | Unknown | No correction identified in bounded checks\par
Human review | None supplied | `human\_reviews=[]`\par
9 Z. Zhang, Maximal parahoric arithmetic transfers, resolutions and modularity, arXiv:2112.11994v3 (https://arxiv.org/abs/2112.11994v3), 9 May 2024. Exact reviewed author manuscript, 106 pages. Associated journal DOI: 10.1215/00127094-2024-0023. W. Zhang, Weil representation and Arithmetic Fundamental Lemma, Ann. of Math. (2) 193 (2021), 863–978, doi:10.4007/annals.2021.193.3.5 (https://doi.org/10.4007/annals.2021.193.3.5). Consulted 116-page typeset PDF from the YMSC academic archive; Appendix A, physical pp. 107–109. P. van Hoften, with an appendix by R. Zhou, Mod p points on Shimura varieties of parahoric level, arXiv:2010.10496v1 (https://arxiv.org/abs/2010.10496v1), 20 October 2020. Consulted setup and connectedness/irreducibility interfaces; no full independent reproduction of its proof.\par
\section*{Supporting checks}
3c242d0b79c457b0d63c5fab12db9cc5
f53f1f1239d585e264f4925ed77c1cd5\par
The exact full digest in the verification records governs the binding. Companion files contain all seven original source bindings, full-read coverage, actual image observations, own exact-calculation scripts and outputs, external versions, build dependencies and logs, and final report-page inspection hashes. The six-lattice sum and unitary matrices use exact arithmetic, not statistical evidence. No author software, author TeX execution or shell escape was used.\par
The root covered its assigned full text and directly checked the central counterexamples; auxiliary roles supplied the remaining full-text scopes and bounded interfaces. Actual inspection does not mean every external theory was reproved. Scientific adjudication remains pending.\par
\end{document}
